\documentclass[10pt,a4paper]{article}
\usepackage[margin=21mm,headheight=14pt]{geometry}
\usepackage[T1]{fontenc}
\usepackage{lmodern,microtype,amsmath,booktabs,array}
\usepackage[dvipsnames]{xcolor}
\usepackage{fancyhdr}
\usepackage[colorlinks=true,linkcolor=MidnightBlue,urlcolor=MidnightBlue]{hyperref}
\pagestyle{fancy}\fancyhf{}
\fancyhead[L]{\small Mihal Dimo / Qwen3.8-27B on H100}
\fancyhead[R]{\small 1 October 2026}
\fancyfoot[C]{\small\thepage}
\setlength{\parindent}{0pt}\setlength{\parskip}{5pt}
\renewcommand{\arraystretch}{1.15}
\newcommand{\MicroRows}{1 & 5.97 & 4.73 & 1.263$\times$ \\
2 & 9.45 & 7.99 & 1.183$\times$ \\
8 & 23.29 & 21.04 & 1.107$\times$ \\
16 & 39.97 & 37.93 & 1.054$\times$ \\
32 & 72.88 & 71.54 & 1.019$\times$ \\
64 & 139.49 & 138.64 & 1.006$\times$ \\
128 & 272.71 & 272.52 & 1.001$\times$ \\}
\newcommand{\NumericalCases}{42}
\newcommand{\NumericalEdges}{6}
\newcommand{\StateError}{8.34e-06}
\newcommand{\OutputError}{0.000977}
\newcommand{\ProfileCalls}{480}
\newcommand{\ProfileGraphs}{10}
\newcommand{\ProfileMean}{66.74}
\newcommand{\EngineRows}{Single / output tok/s & 72.31 & 72.82 & +0.71\% \\
Prefill / input tok/s & 18168.43 & 18136.11 & -0.18\% \\
Decode / output tok/s & 1642.53 & 1641.93 & -0.04\% \\
Mixed / output tok/s & 989.13 & 1002.76 & +1.38\% \\}
\newcommand{\LatencyRows}{Single & 101.69 & 98.39 & 13.15 & 13.08 \\
Prefill & 8159.12 & 8090.34 & --- & --- \\
Decode & 339.67 & 335.08 & 17.02 & 17.04 \\
Mixed & 1909.02 & 1858.83 & 18.40 & 21.48 \\}
\newcommand{\DecodeGain}{-0.04}
\newcommand{\MixedGain}{+1.38}
\newcommand{\TokenAgree}{232}
\newcommand{\ControlAgree}{231}
\newcommand{\MemoryBase}{50.97}
\newcommand{\MemoryCustom}{50.97}

\title{\vspace{-10mm}\textbf{A Custom Gated DeltaNet Decode Kernel}\\[3pt]
\large Register pressure, state bandwidth, and serving tradeoffs on an H100}
\author{Mihal Dimo\\\small\href{https://github.com/mhdimo/qwen38-h100-lab}{github.com/mhdimo/qwen38-h100-lab}}\date{1 October 2026}
\begin{document}
\maketitle\thispagestyle{fancy}\vspace{-6mm}
\begin{abstract}
We implement a CUDA replacement for Qwen3.8-27B's packed Gated DeltaNet (GDN)
single-token recurrence. Four warps divide each state tile, lowering compiler
register usage from 205 per thread in the profiled reference to 40 in the selected
candidate. Three tile sizes are numerically tested and timed. Large-batch gains
are limited by the cost of reading and writing the FP32 recurrent state.
Against an otherwise matched, already-tuned vLLM engine, measured decode-heavy
throughput changes by \DecodeGain{}\% and mixed throughput by \MixedGain{}\%.
These are incremental kernel results, separate from the earlier concurrency
tuning and SiLU--FP8 experiments. Production services remain unchanged.
\end{abstract}

\section{Measured target and pinned environment}
The H100 80GB HBM3 has driver 610.57.04 and a 700 W power limit. The runtime is
vLLM 0.26.0, PyTorch 2.11.0, Triton 3.6.0 and CUDA 13.2.86. We keep native FP8
weights, BF16 activations and FP32 recurrent state. The cached model revision is
\texttt{017b9c7af6b5689d5dd426a76e0bc077eb5ca20a}.

A prior warmed reference capture identified GDN packed recurrence as 20.45\%
of summed GPU kernel time in ten steady decode graphs: 480 calls averaging
66.50~$\mu$s. The installed source hash was rechecked before this pass and matched
the saved reference. These percentages are sums of GPU durations, not exclusive
fractions of wall time. Nsight Compute hardware counters were not collected.

Each of 48 GDN layers uses 48 value heads with a $128\times128$ FP32 state.
Sixteen Q/K heads are shared in groups of three value heads. At batch 32, a single
layer reads and writes about 201.3 MB of state per step. Dividing this by the
reference profile time gives an approximate 3.03 TB/s state-traffic rate; this
is an arithmetic estimate, not a measured DRAM counter. It explains why reducing
register use need not yield a large concurrent-throughput improvement.

\section{Recurrence and CUDA mapping}
For one value head, normalize $q,k$ with epsilon $10^{-6}$, apply the query scale,
and compute
\[
 d=\exp[-\exp(A)\,\operatorname{softplus}(a+\Delta)],\quad
 \beta=\operatorname{FP32}(\operatorname{BF16}(\operatorname{sigmoid}(b))),
\]
\[
 \bar H=dH,\qquad u=\beta(v-\bar Hk),\qquad H'=\bar H+uk^T,
 \qquad o=H'q.
\]
The intermediate BF16 rounding of $\beta$ is preserved explicitly. State updates
stay FP32; the output is BF16. The operation was already fused in the reference.
Our change is the thread/tile mapping and vectorized state access, not a new
fusion of previously separate operators.

{\small Method references:
\href{https://docs.nvidia.com/cuda/cuda-c-best-practices-guide/}{NVIDIA CUDA Best Practices};
\href{https://docs.vllm.ai/en/latest/design/cuda_graphs/}{vLLM CUDA graphs}.
The pinned numerical reference is the installed vLLM v0.26.0 operator; its
upstream source can be obtained from
\href{https://github.com/vllm-project/vllm}{github.com/vllm-project/vllm}.}

\newpage
\section*{Implementation and numerical checks}
The selected CUDA kernel uses four warps per block and four value rows per warp:
a $16\times128$ state tile. Each lane owns four adjacent key dimensions. Aligned
\texttt{float4} loads/stores move the state, and warp shuffles reduce dot products.
There is no shared-memory reduction or block-wide barrier. The compiler reports
40 registers/thread, zero spill loads/stores and zero barriers. The alternative
32- and 64-row tiles use 62 and 104 registers/thread, respectively. The original
Triton kernel uses one warp for a 32-row tile and 205 registers/thread.

The process-local worker preloads the extension before warmup and CUDA graph
capture. It replaces only packed non-speculative decode inside the existing
opaque attention operation. Prefill, full attention, weights and installed
packages are unchanged. Unsupported layouts and dtypes use the original
operator. Positive state slots must be unique and in range; nonpositive padding
slots yield zero output without state modification. Strided input/gate rows,
strided indices and padded state-slot storage are supported.

Each of the three tiles passes \NumericalCases{} matrix cases over three seeds,
batches 0--128, contiguous and padded layouts, BF16/FP32 gate bias, random and
zero initial states. \NumericalEdges{} additional cases cover int32 indices,
unnormalized Q/K, zero Q/K, extreme gates, unaligned-state fallback and FP16
fallback. We compare both installed vLLM and an independent PyTorch recurrence.
Untouched slots must remain exactly unchanged. Predeclared tolerances are
FP32 state rtol $3\times10^{-4}$ / atol $3\times10^{-5}$ and BF16 output rtol
0.016 / atol 0.002; they were not relaxed to obtain a pass.

Each tile also passes 256 CUDA-graph recurrent steps with changing inputs and
independently updated reference state. Maximum trajectory differences are
\StateError{} in state and \OutputError{} in output. This is tolerance-based
equivalence over tested inputs, not bitwise equality or an unrestricted error bound.
Compute Sanitizer 2026.1.1 reports zero memory errors across 48 direct custom-kernel
calls and eight graph replays, including both index/bias types, normalization
modes, padding and an upper-bound guard. Instrumentation is filtered to the
custom GDN kernels; it does not certify the entire vLLM runtime.

\section{Warmed CUDA-graph microbenchmark}
Graphs contain eight operations on independent states. Ten alternating trials
time 30 graph replays each, divided by 240 calls. The table reports median average
time per call for the selected 16-row tile. Loading, compilation and warmup are
excluded. This fixed-input recurrence converges toward steady state; the full
engine comparison separately exercises changing activations.

\begin{center}\small
\begin{tabular}{@{}rrrr@{}}\toprule
Batch & vLLM ($\mu$s) & Custom ($\mu$s) & Speedup \\\midrule
\MicroRows
\bottomrule\end{tabular}
\end{center}
Active state pools are 25.2 MB at batch 1, 201.3 MB at batch 8 and 805.3 MB at
batch 32. Batch 1 can benefit from L2 residency; large batches exceed L2.
Allocated pool sizes include unused sentinel slots and are reported separately
in the raw data. CUDA-event measurements include graph scheduling and submission
gaps. Tile selection uses the exploratory sweep; raw timings for all variants
are retained. Tiny speed differences should not be interpreted as universal wins.

\newpage
\section{Matched full-model experiment}
Both configurations use 64 sequence capacity, 16 GiB cache, a 16,384-token budget,
chunked prefill and \texttt{FULL\_AND\_PIECEWISE} graphs. Only the GDN worker and
its compilation-cache directory differ. The prior custom SiLU kernel is absent
from both, so its gains are not combined with these results.

The custom trace contains \ProfileCalls{} GDN CUDA calls across
\ProfileGraphs{} steady decode graph launches and zero reference packed-decode
calls in that capture. This verifies real engine dispatch. Profiled requests do
not enter the throughput table.
Its mean GDN duration is \ProfileMean{}~$\mu$s versus 66.50~$\mu$s in the earlier
reference capture: essentially unchanged at batch 32 despite fewer registers.

Each configuration completes 400 measured requests: two repetitions of 128/128
input/output tokens at concurrency one, 4096/1 at concurrency 32, 128/128 at
concurrency 32 and 1024/128 at concurrency 32. Seeded prompt tokens and fixed
output lengths match; the harness rejects partial output and prefix-cache hits.

\begin{center}\small
\begin{tabular}{@{}lrrr@{}}\toprule
Workload / throughput & Tuned vLLM & GDN CUDA & Change \\\midrule
\EngineRows
\bottomrule\end{tabular}
\end{center}
\begin{center}\small
\begin{tabular}{@{}lrrrr@{}}\toprule
& \multicolumn{2}{c}{TTFT p95 (ms)} & \multicolumn{2}{c}{TPOT p95 (ms)}\\
Workload & vLLM & Custom & vLLM & Custom\\\midrule
\LatencyRows
\bottomrule\end{tabular}
\end{center}
TTFT includes scheduling and queueing. TPOT is the per-request mean interval
after its first token, not an individual-token tail. A one-token response has
no decode TPOT. Prefill-only throughput acts as a control because this kernel
does not change prefill computation; any difference is run variation.
Maximum end-of-case GPU memory samples are \MemoryBase{} GiB for the reference
and \MemoryCustom{} GiB for the candidate; these are not allocation peaks.
The mixed-workload TPOT p95 regresses from 18.40 to 21.48 ms, despite higher
aggregate throughput. Concurrent decode throughput is effectively flat.

Generated-token hashes agree for \TokenAgree{}/400 candidate/reference requests.
Two unchanged reference runs agree for \ControlAgree{}/400. Both implementations
pass the three text smoke checks. The reference control demonstrates engine
nondeterminism but does not prove all candidate differences harmless. Production
text, image/video and structured-output quality remain unvalidated.

\section{Interpretation and reproducibility}
Lower register pressure produces clear small-batch operation gains, while state
bandwidth limits large-batch gains. Two sequential repetitions, synthetic traffic,
unlocked GPU clocks and scheduling variation limit confidence in small end-to-end
differences. These results do not establish the globally fastest kernel or a
production deployment recommendation. AsyncLLM throughput bypasses the existing
application admission lock and is not production HTTP capacity.
For the stated many-user decode objective, we retain stock GDN as the preferred
baseline: this candidate does not demonstrate an end-to-end decode win. The CUDA
implementation remains an opt-in research artifact for further experiments.

All CUDA source, failed attempts, raw JSON and compiler logs are retained in
this repository; large binary profiler artifacts are kept in a local archive.
The temporary GPU workspace was hash-verified and removed after the study.
\texttt{gdn\_build\_data.py} regenerates tables and source hashes.
\begin{verbatim}
python gdn_pass/run.py --kernel-action gdn-validate
python gdn_pass/run.py --kernel-action gdn-bench
python gdn_pass/run.py --configs batch64_t16384 \
  batch64_t16384_gdn --repetitions 2
\end{verbatim}
\end{document}
